\appendix
\label{chap:appendix}
\chapter{Anhang (WORK IN PROGRESS)}

\section{Quellcode}
\label{app:code}

Im zugehörigen GitHub-Repository finden sich sämtliche Skripte zur Simulation und zum Training der Agenten sowie eine ausführliche README mit Installations- und Anwendungshinweisen:

\begin{itemize}
  \item \href{https://github.com/fynn-madrian}{github.com/fynn-madrian}
\end{itemize}

\section{Zusätzliche Tabellen}
\label{app:tables}

Die folgenden Tabellen führen weitere Hyperparameter auf, die im Hauptteil nur kurz erwähnt wurden:

\subsection{MARL Parameter}
\begin{table}[htbp]
  \centering
  \caption{Modell- und Trainingshyperparameter}
  \label{tab:marl_model_hyperparams}
  \begin{tabular}{@{}lcl@{}}
    \toprule
    Parameter                         & Symbol            & Wert          \\
    \midrule
    Lernrate                          & $lr$              & $1\times10^{-3}$ \\
    Batch-Größe                       & –                 & 32            \\
    PPO-Schritte pro Update           & –                 & 4             \\
    PPO-Clip-Epsilon                  & $\epsilon$        & 0,2           \\
    Diskontfaktor                     & $\gamma$          & 0,95          \\
    GAE-Faktor                        & $\lambda$         & 0,95          \\
    Value-Koeffizient                 & $c_v$             & 0,5           \\
    Entropie-Koeffizient              & $c_e$             & 0,01          \\
    \bottomrule
  \end{tabular}
\end{table}


\subsection{Single Agent Parameter}
\begin{table}[htbp]
  \centering
  \caption{Netzwerarchitektur und Datenhandling}
  \label{tab:network_architecture}
  \begin{tabular}{ll}
    \toprule
    Parameter             & Wert     \\
    \midrule
    Batch-Größe           & 128       \\
    PPO-Schritte pro Update & 3         \\
    Buffer-Größe          & 512      \\
    Sequenzlänge          & 128      \\
    Dropout-Rate          & 0{,}1    \\
    \bottomrule
  \end{tabular}
\end{table}

\begin{table}[htbp]
  \centering
  \caption{RL-Algorithmus und Trainingshyperparameter}
  \label{tab:rl_hyperparams}
  \begin{tabular}{@{}lcc@{}}
    \toprule
    Parameter                            & Symbol                  & Wert       \\
    \midrule
    Diskontfaktor                        & \(\gamma\)              & 0,99       \\
    GAE-Faktor                           & \(\lambda\)             & 0,95       \\
    PPO-Clip                             & \(\epsilon\)            & 0,1        \\
    Entropie-Koeff. (Start)             & \(\beta_{\mathrm{start}}\)   & 0,01       \\
    Entropie-Koeff. (Ende)              & \(\beta_{\mathrm{min}}\)     & 0,001      \\
    Learning Rate (Start)               & \(\alpha_{\mathrm{start}}\)  & \(1\times10^{-5}\)  \\
    Learning Rate (Ende)                & \(\alpha_{\mathrm{min}}\)    & \(1\times10^{-6}\)  \\
    Environment Steps                   & –                       & 5 000 000  \\
    \bottomrule
  \end{tabular}
\end{table}

\section{Animationsclips}
\label{app:animations}

Um die dynamischen Verhaltensweisen der trainierten Modelle zu illustrieren, wurden ausgewählte Durchläufe als Animationsclips exportiert. Die Videodateien sind im Repository unter dem Verzeichnis \texttt{/animations} abrufbar:

\begin{itemize}
\item Link to Gif
\end{itemize}
